In this section we discuss our Facebook Link Recommendation
(LinkR) application, definitions for 
notation in the context of this application, and our evaluation
methodology.

\subsection{Link Recommendation App on Facebook}

To evaluate existing and newly proposed (S)CF methods discussed in
this paper, we created a Facebook application (i.e., a Facebook
``App'') that recommends links to users everyday, where the users may
give their feedback on the links indicating whether they \emph{liked}
it or \emph{disliked} it.  Figure~\ref{fig:linkr_app} shows 
our Facebook LinkR App as it appears to users.

%%%%%%%%%%%%%%%%%%%%%%%%%%%%%%%%%%%%%%%%%%%%%%%%%%%%%%%%%%%%%%%%%%%%
\begin{figure}[t!]
\hspace{-2mm} \subfigure{\includegraphics[scale=0.42]{img/linkr.eps}}
\caption{The Facebook LinkR App showing two link recommendations to a 
user.  One link recommendation is from a friend and includes the
friend's commentary on the link as well as the link description.  The
other link recommendation is from a non-friend and hence only shows
the link description.  Users have the option of liking or disliking each 
recommendation as well as providing explicit feedback.}
\label{fig:linkr_app}
\end{figure}
%%%%%%%%%%%%%%%%%%%%%%%%%%%%%%%%%%%%%%%%%%%%%%%%%%%%%%%%%%%%%%%%%%%%

The functionalities of the LinkR application on a daily basis are as follows:
\begin{enumerate}
\item{Collect data that have been shared by users and their friends on Facebook.}
\item{Initiate retraining of all active (S)CF link recommendation algorithms on the latest collected data.}
\item{Following retraining, recommend three links to the users according to their assigned recommendation algorithm.}
\item{Collect feedback from the users on whether they liked or disliked the recommendations as well as any additional commentary the user wishes to provide.}
\end{enumerate}
Details of the (S)CF link recommendation algorithms and user assignments
will be discussed shortly, but first we cover the specific data
collected by the LinkR App and made available for use by the
recommendation algorithms.

\subsection{Facebook Data Collected}

\label{sec:dataset}

At its peak membership, 108 users had elected to install the Facebook
App developed for this project.  From this user base, we were able to
gather data on 34,860 users and 437,023 links in total by the end of
the evaluation period.

\subsubsection{User Data}

Data that are collected and used to define the user feature vector
$\x$ introduced in Section~\ref{sec:notation} for the LinkR Facebook
App are defined as follows:
\begin{itemize}
\item $\forall \mathit{id}, \, \x_\mathit{id} = \mathit{id} \in \{ 0, 1 \}$: 
every unique Facebook ID (user) recorded in the App was assigned its own
binary indicator in $\x$; all such indicators are enforced to be
mutually exclusive.
\item $\mathit{gender} \in \{ 0 \mbox{ female}, 1 \mbox{ male} \}$.
\item $\mathit{age} \in \mathbb{N}$.
%\item {$\mathit{location}_\mathit{id}$:} an integer ID corresponding
%to the user's specific present location (city and country)
%\item {$\mathit{hometown}_\mathit{id}$:} an integer ID corresponding
%to the user's specific home town (city and country)
%\item $\forall \y \in \mathit{friends}_{\x}, \, \mathit{friend}_{\x,\y} \in \{ 0 \mbox{ not friend}, 1 \mbox{ friend} \}$: 
\end{itemize}
We note that the indicator of friendships for $\x$ is stored in the
$\mathit{friends}_\x$ set defined in Section~\ref{sec:notation} and
used in various previous objective definitions, but not explicitly 
stored in $\x$.
%Although we did include additional binary features representing a user's
%hometown and present location, these features did not appear to prove helpful 
%in initial experimentation and were not used in subsequent trials.

\subsubsection{Link Data}

Data that are collected and used to define the item feature vector
$\y$ introduced in Section~\ref{sec:notation} for the LinkR Facebook
App are defined as follows:
\begin{itemize}
\item $\forall \mathit{id}, \, \y_\mathit{poster} = \mathit{id} \in \{ 0, 1 \}$:  
binary indicator feature for the $\mathit{id}$ of the user who posted the link;
all such binary indicator features are mutually exclusive.
\item $\forall \mathit{id}, \, \y_\mathit{wall} = \mathit{id} \in \{ 0, 1 \}$:
binary indicator feature for the $\mathit{id}$ of the user on whose wall the link was posted; all such binary indicator features are mutually exclusive.
%\item Text description of the link from the user who posted it.
%\item Text link summary from the metatags on the target link webpage.
\item Count of total link ``likes'' on Facebook.
\item Count of total link shares on Facebook.
\item Count of total link comments posted on Facebook.
\end{itemize}

\subsubsection{Joint User and Link Data}

The feature vector $\f_{\x,\y}$ used in Sections~\ref{sec:cbf}
and~\ref{sec:hybrid_scf} for the LinkR Facebook App is defined as the
\emph{concatenation} of $\x$, $\y$ (above) and the following
additional social network information diffusion features:
\begin{itemize}
\item $\forall \z \in \mathit{friends}_\x, \, \mbox{$\z$ liked $\x$} \in \{0, 1\}$: for every friend $\z$ of user $\x$, we have a binary information diffusion  feature indicating whether user $\z$ liked item $\y$ (recall that $\f_{\x,\y}$ is built w.r.t.\ a specific user $\x$ and item $\y$).
\end{itemize}

\subsubsection{Interaction Data}
\label{sec:interactions}

The Facebook social network interactions between users $\x$ and $\z$ that we count
(equally weighted) to define
\emph{\# interactions between $\x$ and $\z$} in Section~\ref{sec:notation} 
are defined as follows:
\begin{enumerate}
\item{Being friends.}
\item{Posting, liking, or commenting on an item (link, photo, video, photo, or message) on a user's wall.}
\item{Being tagged together in the same photo or video.}
\item{Tagging themselves as attending the same school or class, playing sports together, working together for the same company or on the same project.}
\end{enumerate}

\subsection{Live Online Recommendation Trials}

For the recommendations made to the LinkR application users, we select
only links posted in the most recent two weeks that the user has not
already liked (or disliked).  We use only the links from the last two
weeks since an informal user study has indicated a preference for
recent links.  Indeed, older links have a greater chance of being
outdated and are also likely to represent broken links that are no longer
accessible.  We have chosen to recommend three links per day
with the aim of avoiding position bias and information overload that
may occur with longer recommendation lists.
%to the LinkR users and according to the survey done at the end of the
%first trial, three links per day seems to be the generally preferred
%number of daily recommendations.

For the live trials, Facebook users who installed the LinkR
application were \emph{randomly assigned one of four algorithms in
each of two trials}.  Specific algorithms trialed will be discussed in
Section~\ref{sec:EmpResults}.  LinkR users were not informed which
algorithm was assigned to them.  As demonstrated in
Figure~\ref{fig:linkr_app}, we distinguish our recommended links into
two major classes: (1) links that were posted by the LinkR user's
friends, and (2) links that were posted by users other than the LinkR
user's friends.  LinkR users were encouraged to rate each link as like
or dislike.  In turn these ratings became part of the training data
for the recommendation algorithms, and thus were used to improve the
performance of the algorithms over time.  LinkR users were allowed to
provide feedback comments on specific links if they wished; based on
repeated user commentary from the first trial period, we filtered out
non-English links and links lacking any textual descriptions from the
recommendations made on the second trial in order to prevent user
annoyance.

%At the end of the first trial, we conducted a user survey with the
%LinkR users to find out how satisfied they were with the
%recommendations they were getting.
 